\cvsection{Exp\'eriences techniques}

\begin{cventries}

  \cventry
    {DyNaMo, Aix-Marseille Universit\'e}
    {D\'eveloppeur logiciel \& mat\'eriel LabVIEW, doctorant}
    {Marseille, France}
    {janv. 2022 -- actuel}
    {
      \begin{cvitems}
        \item {\textbf{Interdisciplinaire}: collaboration directe avec des chercheurs exp\'erimentaux pour d\'efinir et optimiser les workflows de donn\'ees force-indentation, en traduisant des exigences mat\'erielles complexes en syst\`emes d'acquisition utilisables}
        \item{\textbf{D\'eveloppement logiciel FPGA} (ao\^ut 2023 -- actuel): conception et maintenance de modules personnalis\'es pour syst\`eme HS-AFM, optimisation des transferts DMA pour l'acquisition \`a haute vitesse, gestion de flux \`a haut d\'ebit entre FPGA et h\^ote via PCIe, impl\'ementation d'un buffer DRAM pour la sortie de formes d'onde [LabVIEW]}
        \item{\textbf{NeoLaura - logiciel HS-AFM} (janv. 2019 -- ao\^ut 2022): mise \`a niveau et refonte d'une nouvelle interface utilisateur pour la spectroscopie de force [LabVIEW]}
        \item {\textbf{AFM maison sur Jetson NANO C100} (janv. 2021 -- ao\^ut 2022): mise en place d'un syst\`eme Linux sur Jetson NANO C100 pour un AFM maison avec Red Pitaya connect\'e en LAN pour contr\^ole \`a distance et acquisition [Linux, VHDL]}
        \item {\textbf{OpenFMLab} (janv. 2022 -- actuel): encadrement d'un \`etudiant de master pour la conception d'un scanner XY imprim\'e en 3D avec gains m\'ecaniques pour augmenter la plage}
        \item {\textbf{PyFMlab} (janv. 2022 -- actuel): adaptation et maintenance d'un logiciel Python pour le contr\^ole du syst\`eme AFM, l'acquisition, l'analyse et la visualisation}
      \end{cvitems}
      \vspace{1em}
    }

  \cventry
    {LAI, Aix-Marseille Universit\'e}
    {D\'eveloppeur logiciel \& mat\'eriel LabVIEW, ing\'enieur de recherche}
    {Marseille, France}
    {janv. 2019 -- janv. 2022}
    {
      \begin{cvitems}
        \item{\textbf{Laura - logiciel HS-AFM} (janv. 2019 -- ao\^ut 2022): mise \`a niveau et refonte d'une ancienne interface de spectroscopie de force coupl\'ee au syst\`eme HS-AFM d'Ando (RIBM) [LabVIEW], impl\'ementation de la m\'ethode retract-chirp, analyse de spectre de puissance et optimisation de la sauvegarde des donn\'ees}
        \item {\textbf{Technicien} (janv. 2022 -- actuel): r\'eparation et remplacement de pi\'ezos Z pour scanners HS-AFM d'imagerie}
        \item Analyse de donn\'ees AFM avec Python.
      \end{cvitems}
      \vspace{1em}
    }

  \cventry
    {Integrated Micro-electronics, Inc. (IMI)}
    {D\'eveloppeur logiciel LabVIEW \& ing\'enieur d\'eveloppement de syst\`emes de test}
    {Bi\~nan, Laguna, Philippines}
    {janv. 2016 -- ao\^ut 2018}
    {
      \begin{cvitems}
        \item{\textbf{Int\'egration de capteur pour testeur de tenue en tension AC} (juin -- ao\^ut 2018): support de 2 projets avec d\'eveloppement d'interfaces GUI personnalis\'ees via protocoles GPIB et RS232 pour capteurs et testeurs tiers de tenue en tension (TDK), impl\'ement\'es en LabVIEW avec int\'egration SQL}
        \item{\textbf{D\'eveloppement syst\`eme test HTGB/HTRB/H3TRB pour IGBT} (ao\^ut 2017 -- ao\^ut 2018): conception d'une GUI LabVIEW pour int\'egration de sous-ensembles, impl\'ementation de tests HT/GB et HTTRB pour validation de modules de puissance, ajout de tra\c{c}abilit\'e SQL, conception d'un bo\^itier de test sur mesure}
        \item{\textbf{D\'eveloppement contr\^oleur DUT} (d\'ec. 2017 -- ao\^ut 2018): d\'eveloppement d'un syst\`eme de test pour dispositif IoT Scheduler en d\'ebug et troubleshooting mat\'eriel via CAN, int\'egration de t\^aches IoT avec module Zapier orchestr\'e par JIRA}
        \item{\textbf{D\'eveloppement de drivers mat\'eriels} (janv. 2016 -- mai 2017): mission \`a IMI Jiaxing (Chine) pour soutenir l'\'equipe de d\'eveloppement de syst\`emes de test, d\'eveloppement de drivers oscilloscope sous LabVIEW, m\'ethode LVOOP pour dispatch dynamique des \'equipements en remplacement de la base SQL afin d'assurer l'int\'egrit\'e des donn\'ees}
        \item{\textbf{Testeur optique de lumi\`ere parasite} (nov. 2016 -- mars 2017): revue de l'architecture syst\`eme GUI et logicielle sous LabVIEW pour testeur optique, conception mat\'erielle selon la th\'eorie optique Prindle, conception d'un contr\^oleur PID pour pilotage pneumatique et alignement capteurs du banc Verily}
      \end{cvitems}
      \vspace{1em}
    }

\end{cventries}
